\documentclass[11pt]{article}
\usepackage{times}
\usepackage{fullpage}
\usepackage{url}
\usepackage{hyperref}
\usepackage{fancyhdr}
\usepackage{graphicx}
\usepackage{enumitem}
\usepackage{subcaption}
\usepackage{amsmath, amsfonts, amsthm, fouriernc}
\usepackage{IEEEtrantools}
\usepackage{parskip}

\pagestyle{fancy}
\addtolength{\headheight}{2em}
\addtolength{\headsep}{1.5em}
\lhead{\large{ME 2801}}

\usepackage{sectsty}
\sectionfont{\fontsize{12}{8}\selectfont}

\begin{document}

\begin{center}
  \Large{\bf{Assignment: Block Diagrams and PID Control}}
\end{center}

The textbook exercises, designated with the word ``Nise'', are from the
``Problems'' section at the end of the chapter in \emph{Control System
Engineering} by Norman Nise, 7th edition.

% ---------------------------------------------------------------
\section*{Nise 5.1}

% ---------------------------------------------------------------
\section*{Nise 5.11}

Do this problem both manually and with MATLAB.

\begin{enumerate}
\item By hand, find the closed-loop transfer function using block diagram
  algebra, then calculate the requested metrics using the key equations.
\item Using MATLAB, use \texttt{feedback()} and \texttt{stepinfo()} to obtain
  the same results.
\end{enumerate}

% ---------------------------------------------------------------
\section*{Exercise 3: Proportional Speed Control of the USV}

This exercise connects block diagram algebra to the feedback controllers
we will implement in the USV field labs.
For background on the surge plant model and proportional control,
refer to the \emph{PID Control: From Analysis to Practice} article distributed
in class.

\subsection*{Background}

The surge (forward-speed) dynamics of the USV can be modeled as a first-order
plant:
\[
  G_{\mathrm{surge}}(s) = \frac{5.1}{1.3\,s + 1}
\]
where the input is normalized throttle ($\pm 1$) and the output is forward
speed (m/s).
We close a unity feedback loop with a proportional controller $C(s) = K_p$.

\subsection*{Tasks}

\begin{enumerate}

\item \textbf{Block diagram.}
  Draw the standard unity-feedback block diagram with $G_{\mathrm{surge}}(s)$
  in the forward path and $K_p$ as the controller block.
  Label $R(s)$, $E(s)$, $U(s)$, and $Y(s)$.

\item \textbf{Closed-loop transfer function.}
  Derive $T(s) = Y(s)/R(s)$.  Rewrite in standard first-order form and express
  the closed-loop DC gain $K_{\mathrm{cl}}$ and time constant $\tau_{\mathrm{cl}}$
  as explicit functions of $K_p$.

\item \textbf{Design for rise time.}
  Find the minimum $K_p$ such that the closed-loop rise time satisfies
  $T_r < 1.5$\,s.  Derive $K_{p,\mathrm{min}}$ analytically before evaluating
  numerically.

\item \textbf{MATLAB verification.}
  Use \texttt{feedback()} to build $T(s)$ and \texttt{stepinfo()} to extract
  metrics.  Plot the closed-loop unit-step response for $K_{p,\mathrm{min}}$.

\item \textbf{Report.}
  State the following values from your MATLAB results:
  \begin{itemize}
    \item $K_{p,\mathrm{min}}$
    \item Closed-loop rise time $T_r$
    \item Steady-state error $e_{ss} = 1 - K_{\mathrm{cl}}$
  \end{itemize}
  Briefly explain why $e_{ss} > 0$ for any finite $K_p$ with proportional-only
  control.

\end{enumerate}

% ---------------------------------------------------------------
\section*{Exercise 4: Feedforward Compensation}

This exercise extends Exercise~3 by adding a feedforward path to eliminate
steady-state error without changing the closed-loop dynamics.

\subsection*{Background}

In Exercise~3, proportional-only control always leaves a nonzero steady-state
error.
Increasing $K_p$ arbitrarily is impractical for several reasons:

\begin{itemize}
  \item \textbf{Actuator saturation.}  Normalized throttle is bounded to
        $\pm 1$.  At $t=0^+$, the error is $E(0)=1$, so the initial command is
        $U(0) = K_p$.  Any $K_p > 1$ immediately saturates the throttle.
  \item \textbf{Noise amplification.}  High gain passes sensor noise directly
        to the actuator, causing excessive thruster activity.
  \item \textbf{Reduced stability margins.}  High loop gain reduces gain and
        phase margins, making the closed loop more sensitive to unmodeled
        dynamics.
\end{itemize}

A better strategy is to keep $K_p$ modest and add a feedforward gain $K_{ff}$
to handle the steady-state offset.

\subsection*{Tasks}

\begin{enumerate}

\item \textbf{Feedforward block diagram.}
  A feedforward gain $K_{ff}$ is added from $R(s)$ directly to the plant-input
  summing junction, so the control effort becomes
  \[
    U(s) = K_p\,E(s) + K_{ff}\,R(s).
  \]
  Draw the modified block diagram and label all signals and blocks.

\item \textbf{Closed-loop transfer function with feedforward.}
  Derive $T_{ff}(s) = Y(s)/R(s)$.  Show that choosing
  \[
    K_{ff} = \frac{1}{K_{\mathrm{plant}}}, \qquad K_{\mathrm{plant}} = G_{\mathrm{surge}}(0),
  \]
  gives $T_{ff}(0) = 1$ (zero steady-state error) while leaving
  $\tau_{\mathrm{cl}}$ unchanged.

\item \textbf{MATLAB comparison.}
  Using the same $K_p$ from Exercise~3:
  \begin{enumerate}
    \item Build the feedforward closed-loop \texttt{tf} object and verify with
          \texttt{dcgain()}.
    \item Plot the unit-step responses of the P-only and P+feedforward
          controllers on the same axes.
    \item Report rise time and steady-state value for the feedforward case.
  \end{enumerate}

\item \textbf{Report.}
  State the following values:
  \begin{itemize}
    \item $K_{ff}$ (analytical value)
    \item Closed-loop rise time $T_r$ for P+feedforward
    \item Steady-state error $e_{ss}$ for P+feedforward
  \end{itemize}
  Comment on how $T_r$ compares to the P-only result and explain why.

\end{enumerate}

\end{document}
